\begin{center}
{\Large \bfseries \textcolor{lyred}{第八章 空间解析几何与向量代数}}
\end{center}

\vspace{6pt}
{\large \bfseries \textcolor{lyred}{第8.1 节 向量及其线性运算}}

\vspace{4pt}
{\bfseries \textcolor{lyred}{一.向量的概念}}

{\bfseries \textcolor{lyred}{1.向量:既有大小又有方向的量,记为a, AB}}

\xi \xi \xi \rho 
. 
数学上研究的向量都是\textcolor{lyred}{自由向量}. 
\textcolor{lyred}{2.模}:向量的长度,记为a\rho .模为1的向量称为\textcolor{lyred}{单位向量}. 
模为0 的向量称为\textcolor{lyred}{零向量},记为0
\rho 
,规定它的方向是任意的. 
{\bfseries \textcolor{lyred}{3.共线与共面:若,a b}}

\rho 
\rho 
的方向相同或相反,则称它们\textcolor{lyred}{平行}或\textcolor{lyred}{共线},记为
//
a
b
\rho 
\rho 
. 
三个以上的向量,若平移后可位于同一个平面上,则称它们\textcolor{lyred}{共面}. 
{\bfseries \textcolor{lyred}{4.夹角:设a}}

OA
=
\xi \xi \xi \rho 
\rho 
,b
OB
=
\xi \xi \xi \rho 
\rho 
,则
(0
)
AOB
\varphi 
\varphi 
\pi 
=
\le 
\le 
称为,a b
\rho 
\rho 
的\textcolor{lyred}{夹角},记为\cdot 
( , )
a b
\rho 
\rho 
. 
\cdot 
( , )
a
b
a b
\pi 
\bot
\Leftrightarrow 
=
\rho 
\rho 
\rho 
\rho 
;
\cdot 
//
( , )
a
b
a b
\Leftrightarrow 
=
\rho 
\rho 
\rho 
\rho 
或\pi . 
\vspace{4pt}
{\bfseries \textcolor{lyred}{二.向量的线性运算}}

{\bfseries \textcolor{lyred}{1.向量的加减法}}

\textcolor{lyred}{加法定义}:平行四边形法则或三角形法则. 
\textcolor{lyred}{运算律}.(1)\textcolor{lyred}{交换律}:a
b
b
a
+
=
+
\rho 
\rho 
\rho 
\rho ; 
(2)\textcolor{lyred}{结合律}:(
)
(
)
a
b
c
a
b
c
+
+
=
+
+
\rho 
\rho 
\rho 
\rho 
\rho 
\rho . 
\textcolor{lyred}{减法定义}:
(
)
a
b
a
b
-
=
+-
\rho 
\rho 
\rho 
\rho 
,其中b
-
\rho 
表示b
\rho 
的\textcolor{lyred}{负向量}. 
\textcolor{lyred}{注}.0
a
a
+
=
\rho 
\rho 
\rho ,
a
a
-
=
\rho 
\rho 
\rho 
. 
\textcolor{lyred}{三角不等式}: a
b
a
b
+
\le 
+
\rho 
\rho 
\rho 
\rho 
;当且仅当,a b
\rho 
\rho 
同向时等号成立. 
\textcolor{lyred}{例}.证明:平面上对角线互相平分的四边形为平行四边形. 
证.设对角线交点为M ,则由AB
AM
MB
=
+
\xi \xi \xi \rho 
\xi \xi \xi \xi \rho 
\xi \xi \xi \rho 
, DC
DM
MC
=
+
\xi \xi \xi \rho 
\xi \xi \xi \xi \rho 
\xi \xi \xi \xi \rho 
,得AB
DC
=
\xi \xi \xi \rho 
\xi \xi \xi \rho 
,由 
AD
AM
MD
=
+
\xi \xi \xi \rho 
\xi \xi \xi \xi \rho 
\xi \xi \xi \xi \rho 
, BC
BM
MC
=
+
\xi \xi \xi \rho 
\xi \xi \xi \xi \rho 
\xi \xi \xi \xi \rho 
,得AD
BC
=
\xi \xi \xi \rho 
\xi \xi \xi \rho 
,证毕. 
{\bfseries \textcolor{lyred}{2.向量与数的乘法}}

\textcolor{lyred}{定义}.设\lambda \in \mathbb{R},则\lambda 与a\rho 的\textcolor{lyred}{数乘}a
\lambda \rho 表示这样的一个向量:(1)
a
a
\lambda 
\lambda 
=
\rho 
\rho ; 
(2)
//
a
a
\lambda \rho 
\rho ,当
\lambda >
时a
\lambda \rho 与a\rho 同向,当
\lambda <
时a
\lambda \rho 与a\rho 反向. 
\textcolor{lyred}{注}.0
a
\cdots 
=
\rho 
\rho 
,
\lambda \cdots 
=
\rho 
\rho 
,1 a
a
\cdots 
=
\rho 
\rho ,(
)
1 a
a
-
\cdots 
=-
\rho 
\rho . 
\textcolor{lyred}{运算律}.(1)\textcolor{lyred}{结合律}: (
)
(
)
a
a
\lambda \mu 
\lambda \mu 
=
\rho 
\rho ; 
(2)\textcolor{lyred}{分配律}: (
)
a
b
a
b
\lambda 
\lambda 
\lambda 
+
=
+
\rho 
\rho 
\rho 
\rho 
,(
)a
a
a
\lambda 
\mu 
\lambda 
\mu 
+
=
+
\rho 
\rho 
\rho . 
\textcolor{lyred}{应用}.(1)\textcolor{lyred}{单位化}:若
a \neq 
\rho 
\rho 
,则
ae
a
a
=
\rho 
\rho 
\rho 
是与a\rho 同向的单位向量,称为a\rho 的\textcolor{lyred}{单位化}. 
(2)\textcolor{lyred}{平行的条件}:设
a \neq 
\rho 
\rho 
,则
/ /
!
b
a
\lambda 
\Leftrightarrow \exists 
\in 
\rho 
\rho 
\mathbb{R},使得b
a
\lambda 
=
\rho 
\rho . 
\vspace{4pt}
{\bfseries \textcolor{lyred}{三.空间直角坐标系}}

{\bfseries \textcolor{lyred}{1.坐标系的建立}}

给定点O ,称为\textcolor{lyred}{坐标原点},及三个互相垂直的单位向量, ,
i j k
\rho 
\rho \rho 
,称为\textcolor{lyred}{基本单位向量}, 
它们确定了三根互相垂直的数轴,分别称为, ,
x y z\textcolor{lyred}{ 轴},这样就构成了Oxyz\textcolor{lyred}{ 坐标系}, 
也称为
, , ,
O i j k




\rho 
\rho \rho 
\textcolor{lyred}{系}.习惯上,我们采用\textcolor{lyred}{右手系},即, ,
i j k
\rho 
\rho \rho 
的方向满足\textcolor{lyred}{右手法则}. 
三根坐标轴确定了三个\textcolor{lyred}{坐标平面}.三个坐标平面将空间分成八个\textcolor{lyred}{卦限}. 
{\bfseries \textcolor{lyred}{2.向量的坐标}}

设r
OM
=
\xi \xi \xi \xi \rho 
\rho 
,则以OM 为对角线,以三根坐标轴为棱作长方体,其位于, ,
x y z 轴上的 
三个顶点分别为
,
,
P Q R ,记OP
xi
=
\xi \xi \xi \rho 
\rho ,OQ
yj
=
\xi \xi \xi \rho 
\rho 
,OR
zk
=
\xi \xi \xi \rho 
\rho 
,则 
OM
OP
OQ
OR
xi
yj
zk
=
+
+
=
+
+
\xi \xi \xi \xi \rho 
\xi \xi \xi \rho 
\xi \xi \xi \rho 
\xi \xi \xi \rho 
\rho 
\rho 
\rho 
,称为r\rho 的\textcolor{lyred}{坐标分解式},
,
,
xi yj zk
\rho 
\rho 
\rho 
分别称为r\rho 沿 
三根坐标轴方向的\textcolor{lyred}{分向量},而, ,
x y z 称为它的\textcolor{lyred}{分量},或\textcolor{lyred}{坐标},记
(
)
, ,
r
x y z
=
\rho 
. 
由于
(
)
, ,
M
r
OM
x y z
\leftrightarrow 
=
\leftrightarrow 
\xi \xi \xi \xi \rho 
\rho 
,故, ,
x y z 也称为M 的\textcolor{lyred}{坐标},记
(
)
, ,
M x y z . 
\vspace{4pt}
{\bfseries \textcolor{lyred}{四.利用坐标作向量的线性运算}}

\textcolor{lyred}{定理}.设
(
)
,
,
x
y
z
a
a a a
=
\rho 
,
(
)
,
,
x
y
z
b
b b b
=
\rho 
,\lambda \in \mathbb{R},则 
(1)
(
)
,
,
x
x
y
y
z
z
a
b
a
b a
b a
b
\pm 
=
\pm 
\pm 
\pm 
\rho 
\rho 
;(2)
(
)
,
,
x
y
z
a
a
a
a
\lambda 
\lambda 
\lambda 
\lambda 
=
\rho 
. 
\textcolor{lyred}{推论}.设
a \neq 
\rho 
\rho 
,则
//
y
x
z
x
y
z
b
b
b
b
a
a
a
a
\Leftrightarrow 
=
=
\rho 
\rho 
. 
\textcolor{lyred}{推论}.设r
AB
=
\xi \xi \xi \rho 
\rho 
,若
(
)
,
,
A
x y z
=
,
(
)
,
,
B
x
y
z
=
,则
(
)
,
,
r
x
x y
y z
z
=
-
-
-
\rho 
. 
\textcolor{lyred}{例}.设
(
)
,
,
A x y z
, (
)
,
,
B x
y
z
,
\lambda \neq -,在AB 上求一点M ,使得AM
MB
\lambda 
=
\xi \xi \xi \xi \rho 
\xi \xi \xi \rho 
. 
解.设
(
)
, ,
M x y z ,则
(
)
,
,
AM
x
x y
y z
z
=
-
-
-
\xi \xi \xi \xi \rho 
,
(
)
,
,
MB
x
x y
y z
z
=
-
-
-
\xi \xi \xi \rho 
,故 
(
)
(
)
(
)
,
,
x
x
x
x
x
x
y
y
z
z
AM
MB
y
y
y
y
x
y
z
z
z
z
z
\lambda 
\lambda 
\lambda 
\lambda 
\lambda 
\lambda 
\lambda 
\lambda 
\lambda 
\lambda 
-
=
-

+
+
+

=
\Rightarrow 
-
=
-
\Rightarrow 
=
=
=

+
+
+
-
=
-

\xi \xi \xi \xi \rho 
\xi \xi \xi \rho 
. 
\vspace{4pt}
{\bfseries \textcolor{lyred}{五.向量的模,方向角与投影}}

{\bfseries \textcolor{lyred}{1.向量的模与两点间的距离公式}}

\textcolor{lyred}{定理}.设
(
)
, ,
r
x y z
=
\rho 
,则
r
x
y
z
=
+
+
\rho 
. 
\textcolor{lyred}{推论(距离公式)}.设
(
)
,
,
A x y z
,
(
)
,
,
B x
y
z
,则 
(
)
(
)
(
)
AB
AB
x
x
y
y
z
z
=
=
-
+
-
+
-
\xi \xi \xi \rho 
. 
\textcolor{lyred}{例}.已知两点(
)
4,0,5
A
和
(
)
7,1,3
B
,求与AB
\xi \xi \xi \rho 
平行的单位向量. 
解.
(
)
3,1, 2
AB =
-
\xi \xi \xi \rho 
,
AB =
\xi \xi \xi 
,故所求向量为
(
)
3,1, 2
e =\pm 
-
\rho 
. 
{\bfseries \textcolor{lyred}{2.方向角与方向余弦}}

\textcolor{lyred}{定理}.设
(
)
, ,
r
x y z
=
\rho 
,记
\partial 
( , )
a i
\alpha =
\rho 
\rho 
,
\cdot 
( , )
a j
\beta =
\rho 
\rho 
,
\cdot 
( , )
a k
\gamma =
\rho 
\rho 
,称为r\rho 的\textcolor{lyred}{方向角},则 
cos
x
x
r
x
y
z
\alpha =
=
+
+
\rho 
,
cos
y
y
r
x
y
z
\beta =
=
+
+
\rho 
,
cos
z
z
r
x
y
z
\gamma =
=
+
+
\rho 
. 
\textcolor{lyred}{推论}.(
)
cos ,cos
,cos
r
r
e
r
\alpha 
\beta 
\gamma =
=
\rho 
\rho 
\rho 
;特别地,
cos
cos
cos
\alpha 
\beta 
\gamma 
+
+
=. 
\textcolor{lyred}{注}.称r\rho 的三个方向角的余弦为r\rho 的\textcolor{lyred}{方向余弦}. 
\textcolor{lyred}{例}.已知两点(
)
2,2,
A
和
(
)
1,3,0
B
,求AB
\xi \xi \xi \rho 
的方向角. 
解.
(
)
1,1,
AB =-
-
\xi \xi \xi \rho 
,
AB =
\xi \xi \xi \rho 
,
cos
\alpha =-
,
cos
\beta =
,
cos
\gamma =-
,故 
\pi 
\alpha =
,
\pi 
\beta =
,
\pi 
\gamma =
. 
\textcolor{lyred}{例}.在第一卦限求一点A ,使得OA
\xi \xi \xi \rho 
与x , y 轴的夹角分别为3
\pi , 4
\pi ,且
OA =
\xi \xi \xi \rho 
. 
解.
cos
\alpha =
,
cos
cos
\beta 
\gamma 
=
\Rightarrow 
=
,
(
)
2 1
,
,
3,3 2,3
OA


=
=






\xi \xi \xi \rho 
,故 
(
)
3,3 2,3
A
. 
{\bfseries \textcolor{lyred}{3.向量在轴上的投影}}

\textcolor{lyred}{定义}.设
b \neq 
\rho 
\rho 
,记
\cdot 
Prj
cos( , )
b a
a
a b
=
\rho 
\rho 
\rho 
\rho 
\rho 
,称为a\rho 在b
\rho 
上的\textcolor{lyred}{投影},也记()b
a
\rho 
\rho 
. 
若点O 及单位向量e\rho 确定了u 轴,则a\rho 在u 轴上的\textcolor{lyred}{投影}Prj
Prj
u
e
a
a
=
\rho 
\rho 
\rho ,也记为()u
a\rho 
. 
\textcolor{lyred}{定理}.设
c \neq 
\rho 
\rho 
,则
(
)
Prj
Prj
Prj
c
c
c
a
b
a
b
\lambda 
\mu 
\lambda 
\mu 
+
=
+
\rho 
\rho 
\rho 
\rho 
\rho 
\rho 
\rho 
. 
\textcolor{lyred}{补充练习 }
1.求以a
i
j
=
+
\rho 
\rho 
\rho 
,
b
j
k
=-
+
\rho 
\rho 
\rho 
为边的平行四边形的对角线长度. 
解.
(
)
1, 1,1
a
b
+
=
-
=
\rho 
\rho 
,
(
)
1,3, 1
a
b
-
=
-
=
\rho 
\rho 
. 
2.设
(
)
7, 4, 4
a =
-
-
\rho 
,
(
)
2, 1,2
b =-
-
\rho 
,求c\rho ,使
3 42
c =
\rho 
,且c\rho 位于a\rho ,b
\rho 
之间夹角的 
角平分线上. 
解.
a =
\rho 
,
b =
\rho 
,故
(
)
(
)
3 42
7 1, 7,2
c
a
b
=\pm 
+
=\pm 
-
\rho 
\rho 
\rho 
. 
3.设, ,
a b c
\rho 
\rho 
\rho 均为非零向量,两两不共线,但a
b
+
\rho 
\rho 
与c\rho 共线,b
c
+
\rho 
\rho 与a\rho 共线,证明: 
a
b
c
+
+
=
\rho 
\rho 
\rho 
\rho 
. 
证. a
b
c
\lambda 
+
=
\rho 
\rho 
\rho ,b
c
a
\mu 
+
=
\rho 
\rho 
\rho ,故
a
b
c
c
c
a
a
\lambda 
\mu 
\lambda 
\mu 
+
+
=
+
=
+
\Rightarrow 
=
=-
\rho 
\rho 
\rho 
\rho 
\rho 
\rho 
\rho 
,证毕. 
\vspace{6pt}
{\large \bfseries \textcolor{lyred}{第8.2 节 数量积 向量积 混合积}}

\vspace{4pt}
{\bfseries \textcolor{lyred}{一.两向量的数量积}}

\textcolor{lyred}{定义}.
\cdot 
os( , )
Prjb
a b
a b
a b
b
a
\cdots 
=
=
\rho 
\rho 
\rho 
\rho 
\rho 
\rho 
\rho 
\rho 
\rho ,称为a\rho ,b
\rho 
的\textcolor{lyred}{数量积(点积)}. 
\textcolor{lyred}{几何意义}.
\cdot 
cos( , )
a b
a b
a b
\cdots 
=
\rho 
\rho 
\rho 
\rho 
\rho 
\rho 
, Prjb
a b
a
b
\cdots 
=
\rho 
\rho 
\rho 
\rho 
\rho . 
\textcolor{lyred}{物理意义}.设物体在常力F
\rho 
作用下,沿直线从点A 移动到点B ,则F
\rho 
所作的功为 
cos
cos
W
F
AB
F AB
F AB
\theta 
\theta 
=
\cdots 
=
=
\cdots 
\xi \xi \xi \rho 
\xi \xi \xi \rho 
\xi \xi \xi \rho 
\rho 
\rho 
\rho 
. 
\textcolor{lyred}{性质}.(1)
a a
a
\cdots 
=
\rho \rho 
\rho ;(2)
a
b
a b
\bot
\Leftrightarrow 
\cdots 
=
\rho 
\rho 
\rho 
\rho 
. 
\textcolor{lyred}{运算律}.(1)\textcolor{lyred}{交换律}:a b
b a
\cdots 
=
\cdots 
\rho 
\rho 
\rho 
\rho ; 
(2)\textcolor{lyred}{分配律}:(
)
a
b
c
a c
a c
+
\cdots 
=
\cdots +
\cdots 
\rho 
\rho 
\rho 
\rho \rho 
\rho \rho ,
(
)
a
b
c
a b
a c
\cdots 
+
=
\cdots 
+
\cdots 
\rho 
\rho 
\rho 
\rho 
\rho 
\rho \rho ; 
(3)\textcolor{lyred}{结合律}:(
)
(
)
(
)
a
b
a
b
a b
\lambda 
\lambda 
\lambda 
\cdots 
=
\cdots 
=
\cdots 
\rho 
\rho 
\rho 
\rho 
\rho 
\rho 
. 
\textcolor{lyred}{注}.(
)(
)
a
b
a
b
a
a b
b
+
\cdots 
+
=
+
\cdots 
+
\rho 
\rho 
\rho 
\rho 
\rho 
\rho 
\rho 
\rho 
,(
)(
)
a
b
a
b
a
b
+
\cdots 
-
=
-
\rho 
\rho 
\rho 
\rho 
\rho 
\rho 
. 
\textcolor{lyred}{定理(坐标表达式)}.设
(
)
,
,
x
y
z
a
a a
a
=
\rho 
,
(
)
,
,
x
y
z
b
b b b
=
\rho 
,则
x
x
y
y
z
z
a b
a b
a b
a b
\cdots 
=
+
+
\rho 
\rho 
. 
\textcolor{lyred}{推论}.
\cdot 
cos( , )
x
x
y
y
z
z
x
y
z
x
y
z
a b
a b
a b
a b
a
a
a
b
b
b
+
+
=
+
+
+
+
\rho 
\rho 
;
Prj
x
x
y
y
z
z
b
x
y
z
a b
a b
a b
a
b
b
b
+
+
=
+
+
\rho \rho 
. 
\textcolor{lyred}{例}.证明三角形的余弦定理. 
证.
|
|
(
) (
) |
|
|
|
a
BC
BC BC
AC
AB
AC
AB
AC
AB
AC AB
=
=
\cdots 
=
-
\cdots 
-
=
+
-
\cdots 
=
\xi \xi \xi \rho 
\xi \xi \xi \rho \xi \xi \xi \rho 
\xi \xi \xi \rho 
\xi \xi \xi \rho 
\xi \xi \xi \rho 
\xi \xi \xi \rho 
\xi \xi \xi \rho 
\xi \xi \xi \rho 
\xi \xi \xi \rho \xi \xi \xi \rho 
cos
cos
AC
AB
AC AB
A
b
c
bc
A
+
-
=
+
-

\xi \xi \xi \rho 
\xi \xi \xi \rho 
\xi \xi \xi \rho \xi \xi \xi \rho 
,证毕. 
\textcolor{lyred}{例}.已知三点
(
)
1,1,1
M
,
(
)
2,2,1
A
,
(
)
2,1,2
B
,求
AMB

. 
解.
(
)
1,1,0
MA =
\xi \xi \xi \rho 
,
(
)
1,0,1
MB =
\xi \xi \xi \rho 
,故
cos
MA MB
AMB
AMB
MA MB
\pi 
\cdots 

=
=
\Rightarrow 
=
\xi \xi \xi \rho \xi \xi \xi \rho 
\xi \xi \xi \rho \xi \xi \xi \rho 
. 
\vspace{4pt}
{\bfseries \textcolor{lyred}{二.两向量的向量积}}

\textcolor{lyred}{定义}.c
a
b
=
\times 
\rho 
\rho 
\rho 
是这样一个向量:它的模
\cdot 
sin( , )
c
a b
a b
=
\rho 
\rho 
\rho 
\rho 
\rho 
,垂直于a\rho ,b
\rho 
所在平面, 
并且满足\textcolor{lyred}{右手法则},称c\rho 为a\rho ,b
\rho 
的\textcolor{lyred}{向量积(叉积)}. 
\textcolor{lyred}{几何意义}. a b
\times 
\rho 
\rho 
等于以a\rho ,b
\rho 
为边的平行四边形的面积. 
\textcolor{lyred}{性质}.(1)
a
a
\times 
=
\rho 
\rho 
\rho 
;(2)
//
a
b
a
b
\Leftrightarrow 
\times 
=
\rho 
\rho 
\rho 
\rho 
\rho 
. 
\textcolor{lyred}{运算律}.(1)\textcolor{lyred}{反交换律}:a
b
b
a
\times 
=-\times 
\rho 
\rho 
\rho 
\rho ; 
(2)\textcolor{lyred}{分配律}:(
)
a
b
c
a
c
b
c
+
\times 
=
\times 
+
\times 
\rho 
\rho 
\rho 
\rho 
\rho 
\rho 
\rho ,
(
)
a
b
c
a b
a
c
\times 
+
=
\times 
+
\times 
\rho 
\rho 
\rho 
\rho 
\rho 
\rho 
\rho ; 
(3)\textcolor{lyred}{结合律}:(
)
(
)
(
)
a
b
a
b
a
b
\lambda 
\lambda 
\lambda 
\times 
=
\times 
=
\times 
\rho 
\rho 
\rho 
\rho 
\rho 
\rho 
. 
\textcolor{lyred}{注}.(
)(
)
a
b
a
b
a
b
+
\times 
-
=-
\times 
\rho 
\rho 
\rho 
\rho 
\rho 
\rho 
. 
\textcolor{lyred}{定理(坐标表达式)}.设
(
)
,
,
x
y
z
a
a a a
=
\rho 
,
(
)
,
,
x
y
z
b
b b b
=
\rho 
,则 
(
)
,
,
y
z
z
y
z
x
x
z
x
y
y
x
x
y
z
x
y
z
i
j
k
a
b
a b
a b a b
a b a b
a b
a
a
a
b
b
b
\times 
=
-
-
-
=
\rho 
\rho 
\rho 
\rho 
\rho 
. 
\textcolor{lyred}{例}.设, ,
a b c \neq 
\rho 
\rho 
\rho 
\rho 
,证明:b
c
a b
a c
=
\Leftrightarrow 
\cdots 
=
\cdots 
\rho 
\rho 
\rho 
\rho 
\rho \rho , a
b
a
c
\times 
=
\times 
\rho 
\rho 
\rho 
\rho . 
证.必要性:显然.充分性:
(
)
a
b
c
\cdots 
-
=
\rho 
\rho 
\rho 
,
(
)
a
b
c
\times 
-
=
\rho 
\rho 
\rho 
,即b
c
-
\rho 
\rho 与a\rho 垂直又平行, 
故
b
c
-
=
\rho 
\rho 
\rho 
,即得,证毕. 
\textcolor{lyred}{例}.设
(2,1, 1)
a =
-
\rho 
,
(1, 1,2)
b =
-
,求与a\rho ,b
\rho 
均垂直的单位向量e\rho . 
解.
(
)
1, 5, 3
i
j
k
a
b
\times 
=
-=
-
-
-
\rho 
\rho 
\rho 
\rho 
\rho 
,故
(
)
1, 5, 3
e =\pm 
-
-
\rho 
. 
\textcolor{lyred}{例}.已知三点
(
)
1,2,3
A
, (
)
3,4,5
B
,
(
)
2,4,7
C
,求三角形ABC 的面积. 
解.
(
)(
)
(
)
2,2,2
1,2,4
4, 6,2
S
AB
AC
=
\times 
=
\times 
=
-
=
=
\xi \xi \xi \rho 
\xi \xi \xi \rho 
. 
\vspace{4pt}
{\bfseries \textcolor{lyred}{三.向量的混合积}}

\textcolor{lyred}{定义}.称(
)
a
b
c
\times 
\cdots 
\rho 
\rho 
\rho 为a\rho ,b
\rho 
,c\rho 的\textcolor{lyred}{混合积},记为
, ,
a b c




\rho 
\rho 
\rho . 
\textcolor{lyred}{定理(坐标表达式)}.设
(
)
,
,
x
y
z
a
a a
a
=
\rho 
,
(
)
,
,
x
y
z
b
b b b
=
\rho 
,
(
)
,
,
x
y
z
c
c c c
=
\rho 
,则 
, ,
x
y
z
x
y
z
x
y
z
a
a
a
a b c
b
b
b
c
c
c

=


\rho 
\rho 
\rho 
. 
\textcolor{lyred}{推论}.
, ,
, ,
, ,
, ,
a b c
a c b
b a c
c b a








=-
=-
=-








\rho 
\rho 
\rho 
\rho 
\rho 
\rho 
\rho \rho 
\rho \rho 
\rho 
\rho , 
, ,
, ,
, ,
a a b
a b a
a b b






=
=
=






\rho 
\rho 
\rho \rho 
\rho \rho 
\rho 
\rho 
\rho 
,
, ,
, ,
, ,
a b c
b c a
c a b






=
=






\rho 
\rho 
\rho 
\rho 
\rho 
\rho \rho 
\rho \rho 
. 
\textcolor{lyred}{几何意义}.以a\rho ,b
\rho 
,c\rho 为棱的平行六面体的体积
, ,
V
a b c


=

\rho 
\rho 
\rho 
; 
以a\rho ,b
\rho 
,c\rho 为棱的四面体的体积
, ,
V
a b c


=


\rho 
\rho 
\rho 
. 
\textcolor{lyred}{推论}. a\rho ,b
\rho 
,c\rho 共面
, ,
x
y
z
x
y
z
x
y
z
a
a
a
a b c
b
b
b
c
c
c


\Leftrightarrow 
=
\Leftrightarrow 
=


\rho 
\rho 
\rho 
. 
\textcolor{lyred}{推论}.以
(
)
,
,
A x y z
,
(
)
,
,
B x
y
z
, (
)
,
,
C x y z
,
(
)
,
,
D x y z
为顶点的四面体体积 
ABCD
V
x
x
y
y
z
z
x
x
y
y
z
z
x
x
y
y
z
z
-
-
-
=
-
-
-
-
-
-
. 
\textcolor{lyred}{例}.求
(
)
1,2,0
A
,
(
)
2,3,1
B
,
(
)
4,2,2
C
确定的平面\Pi 的方程. 
解.
(
)
, ,
M x y z
\forall 
,
, ,
,
M
A B C M
\in \Pi \Leftrightarrow 
共面
2 1
1 0
4 1
x
y
z
-
-
-
\Leftrightarrow 
-
-
-
=
-
-
-
,即 
x
y
z
+
-
-
=
. 
\textcolor{lyred}{补充练习 }
1.设|
| 2
a =
\rho 
,|
|
b =
\rho 
,|
| 2
a
b
\times 
=
\rho 
\rho 
,且a\rho 与b
\rho 
的夹角\theta 为钝角,求a b
\cdots 
\rho 
\rho 
. 
解.
sin
sin
cos
a b
\theta 
\theta 
\theta 
=
\Rightarrow 
=
\Rightarrow 
=-
\rho 
\rho 
,故
cos
a b
a b
\theta 
\cdots 
=
=-
\rho 
\rho 
\rho 
\rho 
. 
2.设
a \neq 
\rho 
\rho 
,
b \neq 
\rho 
\rho 
,并且
a
b
a
b
+
\bot
-
\rho 
\rho 
\rho 
\rho 
,
a
b
a
b
-
\bot
-
\rho 
\rho 
\rho 
\rho 
,求\cdot 
( , )
a b
\rho 
\rho 
. 
解. (
)(
)
(
)(
)
a
b
a
b
a
a b
b
a b
a
b
a
b


+
\cdots 
-
=
=
\cdots 


\Rightarrow 


=
\cdots 
-
\cdots 
-
=



\rho 
\rho 
\rho 
\rho 
\rho 
\rho 
\rho 
\rho 
\rho 
\rho 
\rho 
\rho 
\rho 
\rho 
,故
\cdot 
\cdot 
cos( , )
( , )
a b
a b
a b
a b
\pi 
\cdots 
=
=
\Rightarrow 
=
\rho 
\rho 
\rho 
\rho 
\rho 
\rho 
\rho 
\rho 
. 
3.设
b =
\rho 
, \cdot 
( , )
a b
\pi 
=
\rho 
\rho 
,求
lim
x
a
xb
a
x
\to 
+
-
\rho 
\rho 
\rho 
. 
解.
(
)
(
)
lim
lim
lim
x
x
x
a
xb
a
a
xb
a
a
xa b
x b
a
x
x a
xb
a
x a
xb
a
\to 
\to 
\to 
+
-
+
-
+
\cdots 
+
-
=
=
=
+
+
+
+
\rho 
\rho 
\rho 
\rho 
\rho 
\rho 
\rho 
\rho 
\rho 
\rho 
\rho 
\rho 
\rho 
\rho 
\rho 
\rho 
\rho 
\cdot 
lim
cos( , )
x
a b
x b
a b
b
a b
a
a
xb
a
\to 
\cdots 
+
\cdots 
=
=
=
+
+
\rho 
\rho 
\rho 
\rho 
\rho 
\rho 
\rho 
\rho 
\rho 
\rho 
\rho 
\rho 
. 
4.设
(
)
2, 3,1
a =
-
\rho 
,
(
)
1, 2,3
b =
-
\rho 
,
(
)
2,1,2
c =
\rho 
,若r
a
\bot
\rho 
\rho ,r
b
\bot
\rho 
\rho 
,Prj
c r =
\rho \rho 
,求r\rho . 
解.设
(
)
, ,
r
x y z
=
\rho 
,则
x
y
z
x
x
y
z
y
x
y
z
z
-
+
=
=




-
+
=
\Rightarrow 
=




+
+
=
=


,故
(
)
14,10,2
r =
\rho 
. 
5.已知(
)
a
b
c
\times 
\cdots 
=
\rho 
\rho 
\rho 
,求(
)(
)
(
)
a
b
b
c
c
a


+
\times 
+
\cdots 
+


\rho 
\rho 
\rho 
\rho 
\rho 
\rho . 
解. (
)(
)
(
)(
)
(
)
(
)
a
b
b
c
c
a
a
b
c
b
c
a
a
b
c


+
\times 
+
\cdots 
+
=
\times 
\cdots 
+
\times 
\cdots 
=
\times 
\cdots 
=


\rho 
\rho 
\rho 
\rho 
\rho 
\rho 
\rho 
\rho 
\rho 
\rho 
\rho 
\rho 
\rho 
\rho 
\rho 
. 
6.设
a =
\rho 
,
b =
\rho 
, \cdot 
( , )
a b
\pi 
=
\rho 
\rho 
,求(1)
(
)
Prja b a
b
-
+
\rho 
\rho 
\rho 
\rho 
;(2)
\cdot 
cos(
,
)
a
b a
b
+
-
\rho 
\rho 
\rho 
\rho 
. 
解.
(
)(
)
3 1
3 1 cos
a
b
a
b
a
b
\pi 
+
=
+
\cdots 
+
=
++
\cdots 
\cdots \cdots 
=
\rho 
\rho 
\rho 
\rho 
\rho 
\rho 
, 
(
)(
)
3 1
3 1 cos
a
b
a
b
a
b
\pi 
-
=
-
\cdots 
-
=
+-\cdots 
\cdots \cdots 
=
\rho 
\rho 
\rho 
\rho 
\rho 
\rho 
, 
(
)(
)
3 1
a
b
a
b
+
\cdots 
-
=
-=
\rho 
\rho 
\rho 
\rho 
; 
(1)
(
)
(
)(
)
Prj
a b
a
b
a
b
a
b
a
b
-
+
\cdots 
-
+
=
=
=
-
\rho 
\rho 
\rho 
\rho 
\rho 
\rho 
\rho 
\rho 
\rho 
\rho 
; 
(2)
\cdot 
(
)(
)
2 7
cos(
,
)
7 1
a
b
a
b
a
b a
b
a
b a
b
+
\cdots 
-
+
-
=
=
=
\cdots 
+
-
\rho 
\rho 
\rho 
\rho 
\rho 
\rho 
\rho 
\rho 
\rho 
\rho 
\rho 
\rho 
. 
7.设
a =
\rho 
,
b =
\rho 
,且
a
b
a
b
-
\bot
-
\rho 
\rho 
\rho 
\rho 
,求以
a
b
+
\rho 
\rho 
和2a
b
+
\rho 
\rho 
为边的三角形面积. 
解.
(
)(
)
\cdot 
\cdot 
sin( , )
sin( , )
A
a
b
a
b
a
b
a b
a b
a b
=
+
\times 
+
=
\times 
=
=
\rho 
\rho 
\rho 
\rho 
\rho 
\rho 
\rho 
\rho 
\rho 
\rho 
\rho 
\rho 
, 
(
)(
)
a
b
a
b
a
a b
b
a b
-
\cdots 
-
=
\Rightarrow 
-
\cdots 
+
=
\Rightarrow 
\cdots 
=
\rho 
\rho 
\rho 
\rho 
\rho 
\rho 
\rho 
\rho 
\rho 
\rho 
,故 
\cdot 
\cdot 
cos( , )
sin( , )
a b
a b
a b
a b
\cdots 
=
=
\Rightarrow 
=
\rho 
\rho 
\rho 
\rho 
\rho 
\rho 
\rho 
\rho 
,故
A =
. 
8.求位于
(
)
2,1,2
a =
\rho 
,
(
)
1, 2,2
b =
-
\rho 
所在平面,与a\rho 垂直的单位向量. 
解.设
(
)
, ,
e
x y z
=
\rho 
,则
x
y
z
+
+
=, 2
x
y
z
+
+
=
, 2
x
y
z
=
-
,解得 
(
)
1, 22,10
e =\pm 
-
\rho 
. 
\vspace{6pt}
{\large \bfseries \textcolor{lyred}{第8.3 节 平面及其方程}}

\vspace{4pt}
{\bfseries \textcolor{lyred}{一.曲面方程与空间曲线方程的概念}}

\textcolor{lyred}{定义}.设有曲面S 和方程
(
)
, ,
F x y z =
,若(
)
(
)
, ,
, ,
x y z
S
F x y z
\in 
\Leftrightarrow 
=
,则称方程 
(
)
, ,
F x y z =
为S 的\textcolor{lyred}{方程},而S 称为方程
(
)
, ,
F x y z =
的\textcolor{lyred}{图形}. 
\textcolor{lyred}{定义}.设有空间曲线\Gamma 和方程组
(
)
(
)
, ,
, ,
F x y z
G x y z
=

=

,若(
)
(
)
(
)
, ,
, ,
, ,
F x y z
x y z
G x y z
=

\in \Gamma \Leftrightarrow 
=

, 
则此方程组称为\Gamma 的\textcolor{lyred}{方程},而\Gamma 称为方程组的\textcolor{lyred}{图形}. 
\vspace{4pt}
{\bfseries \textcolor{lyred}{二.平面的点法式方程}}

\textcolor{lyred}{定理}.过点
(
)
,
,
M
x
y z
,以非零向量
(
)
, ,
n
A B C
=
\rho 
为法向量的平面\Pi 是唯一的, 
它的方程为(
)
(
)
(
)
A x
x
B y
y
C z
z
-
+
-
+
-
=
,称为该平面的\textcolor{lyred}{点法式方程}. 
\textcolor{lyred}{注}.若平面方程为
Ax
By
Cz
D
+
+
+
=
,则
(
)
, ,
n
A B C
=
\rho 
为其法向量. 
\textcolor{lyred}{例}.求过点(
)
2, 3,0
-
,且平行于平面
x
y
z
-
+
+=
的平面方程. 
解.
(
)
1, 2,3
n =
-
\rho 
,故所求方程为(
)
(
)
x
y
z
-
-
+
+
=
,即
x
y
z
-
+
-
=
. 
\textcolor{lyred}{例}.求过点(
)
1,1,1 ,垂直于平面
x
y
z
-
+
=
,3
x
y
z
+
-
+
=
的平面方程. 
解.
(
)(
)
(
)
(
)
1, 1,1
3,2, 12
10,15,5
5 2,3,1
n
n
n
=
\times 
=
-
\times 
-
=
=
\rho 
\rho 
\rho 
,故 
(
)
(
)
(
)
x
y
z
\cdots 
-
+\cdots 
-
+\cdots 
-
=
,即2
x
y
z
+
+
-
=
. 
\textcolor{lyred}{例}.求过三点
(
)
1 2, 1,4
M
-
,
(
)
1,3, 2
M
-
-
,
(
)
3 0,2,3
M
的平面方程. 
解.取
(
)(
)
(
)
3,4, 6
2,3, 1
14,9, 1
n
M M
M M
=
\times 
=-
-
\times -
-
=
-
\xi \xi \xi \xi \xi \xi \rho 
\xi \xi \xi \xi \xi \xi \rho 
\rho 
,故 
(
)
(
)(
)
x
y
z
\cdots 
-
+\cdots 
+
-
-
=
,即14
x
y
z
+
-
-
=
. 
\textcolor{lyred}{注}.一般地,过点
(
)
,
,
M
x
y z
,与不共线的向量a\rho ,b
\rho 
均平行的平面方程为 
0, ,
x
y
z
x
y
z
x x
y y
z z
MM a b
a
a
a
b
b
b
-
-
-

=\Leftrightarrow 
=


\xi \xi \xi \xi \xi \rho 
\rho 
\rho 
;过三点
,
,
M M
M 的平面方程为
,
,
x
x
y
y
z
z
MM M M M M
x
x
y
y
z
z
x
x
y
y
z
z
-
-
-

=
\Leftrightarrow 
-
-
-
=


-
-
-
\xi \xi \xi \xi \xi \rho \xi \xi \xi \xi \xi \xi \rho \xi \xi \xi \xi \xi \xi \rho 
. 
\vspace{4pt}
{\bfseries \textcolor{lyred}{三.平面的一般方程}}

三元一次方程(线性方程)
Ax
By
Cz
D
+
+
+
=
称为平面的\textcolor{lyred}{一般方程},其中 
(1)若
D =
,则平面过原点; 
(2)若
A =
,则平面平行于x 轴,因为
(
)
0, ,
n
B C
=
\rho 
垂直于x 轴;若
B =
,则平面 
平行于y 轴;若
C =
,则平面平行于z 轴. 
(3)若
A
B
=
=
,则平面平行于,x y 轴,即平行于xOy 面. 
\textcolor{lyred}{例}.求过点
(
)
4, 3, 1
M
-
-
与x 轴的平面\Pi 的方程. 
解一.设
:
By
Cz
\Pi 
+
=
,由(
)
4, 3, 1
-
-
\in \Pi ,得
(
)
(
)
B
C
C
B
-
+
-
=
\Rightarrow 
=-
,故 
:
y
z
\Pi 
-
=
. 
解二.
(
)
4, 3, 1 //
OM =
-
-
\Pi 
\xi \xi \xi \xi \rho 
,
(
)
1,0,0 //
i =
\Pi 
\rho 
,可取
(
)
0,1, 3
n
i
OM
=\times 
=
-
\xi \xi \xi \xi \rho 
\rho 
\rho 
,故 
(
)(
)
(
)
x
y
z
\cdots 
-
+
+
-\cdots 
+
=
,即
:
y
z
\Pi 
-
=
. 
\textcolor{lyred}{例}.求过点(
)
1,1,1 ,并且垂直于平面
x
y
z
-
+
=
,3
x
y
z
+
-
+
=
的平面方程. 
解一.设
Ax
By
Cz
D
+
+
+
=
,则
A
B
C
D
A
B
C
A
B
C
+
+
+
=


-
+
=


+
-
=

,解得 
D
A =-
,
D
B =-
,
D
C =-
,故2
x
y
z
+
+
-
=
. 
解二.取
(
)(
)
(
)
(
)
1, 1,1
3,2, 12
10,15,5
5 2,3,1
n =
-
\times 
-
=
=
\rho 
,故 
(
)
(
)
(
)
x
y
z
\cdots 
-
+\cdots 
-
+\cdots 
-
=
,即2
x
y
z
+
+
-
=
. 
\textcolor{lyred}{例}.求过两点(
)
1,1,1
A
,
(
)
0,1, 1
B
-
,且垂直于平面
x
y
z
+
+
=
的平面方程. 
解一.设
Ax
By
Cz
D
+
+
+
=
,则 
A
B
C
D
A
C
B
C
D
B
C
A
B
C
D
+
+
+
=
=-




-
+
=
\Rightarrow 
=




+
+
=
=


,即2
x
y
z
-
+
+
=
. 
解二.取
(
)(
)
(
)
1,0, 2
1,1,1
2, 1, 1
n
AB
n
=
\times 
=-
-
\times 
=
--
\xi \xi \xi \rho 
\rho 
\rho 
,于是 
(
)(
)(
)
x
y
z
-
-
-
-
-
=
,即2
x
y
z
-
-
=
. 
\vspace{4pt}
{\bfseries \textcolor{lyred}{四.平面的截距式方程}}

\textcolor{lyred}{定理}.在三根坐标轴上的截距分别为a ,b ,c 的平面方程为
x
y
z
a
b
c
+
+
=. 
\textcolor{lyred}{例}.求在y 轴和z 轴上截距分别为30和10,并且与
(
)
2,1,3
r =
\rho 
平行的平面方程. 
解.设
x
y
z
a +
+
=,则
(
)
,
,
2,1,3
30 10
a
a
a

\bot
\Rightarrow 
+
+
=
\Rightarrow 
=-




,故 
x
y
z
+
+
=
-
,即5
x
y
z
-
-
+
=
. 
\textcolor{lyred}{例}.求平行于平面2
x
y
z
+
+
+
=
,并且与三个坐标面构成的四面体体积为1的 
平面方程. 
解.设2
x
y
z
D
+
+
+
=
,即
x
y
z
D
D
D
+
+
=
-
-
-
,于是
D
-
=,得 
3 24
D =\pm 
,故
x
y
z
+
+
\pm 
=
. 
\vspace{4pt}
{\bfseries \textcolor{lyred}{五.两平面的夹角}}

\textcolor{lyred}{定义}.两平面的法向量的夹角
\pi 
\theta 
\theta 


\le 
\le 




称为两平面的\textcolor{lyred}{夹角},即 
若
:
A x
B y
C z
D
\Pi 
+
+
+
=
,
:
A x
B y
C z
D
\Pi 
+
+
+
=
,则两者的夹角为 
arccos
A A
B B
C C
A
B
C
A
B
C
\theta 
+
+
=
+
+
\cdots 
+
+
. 
\textcolor{lyred}{例}.求通过x 轴,且与平面
x
y
-
=
成3
\pi 夹角的平面方程. 
解.设
By
Cz
+
=
,由
(
)(
)
0, ,
1, 1,0
cos 3
1 1
B C
B
B
C
B
C
\pi 
\cdots 
-
=
\Rightarrow 
=
+
+
+
\cdots 
++
,得B
C
=\pm 
,故 
所求平面方程为
y
z
\pm 
=
. 
\vspace{4pt}
{\bfseries \textcolor{lyred}{六.点到平面的距离}}

\textcolor{lyred}{定理}.点
(
)
,
,
P x
y z
到平面
:
Ax
By
Cz
D
\Pi 
+
+
+
=
的距离为 
C
B
A
D
Cz
By
Ax
d
+
+
+
+
+
=
. 
\textcolor{lyred}{推论}.两平面
:
Ax
By
Cz
D
\Pi 
+
+
+
=
,
:
Ax
By
Cz
D
\Pi 
+
+
+
=
之间的距离 
D
D
d
A
B
C
-
=
+
+
. 
\textcolor{lyred}{例}.求平面
x
y
z
+
+
=与坐标面构成四面体的内切球球心坐标. 
解.设球心为(
)
,
,
x
y z
,则由它到四面体的四个面距离相等,即 
(
)
1 3
x
y
z
x
y
z
z
z
-
+
+
=
=
=
\Rightarrow 
=-
+
+
,故
x
y
z
=
=
=
+
. 
\textcolor{lyred}{例}.求两平面
x
y
z
+
+
+=
,2
x
y
z
+
+
-
=
的角平分面方程. 
解.
(
)
, ,
1 4
9 1
x
y
z
x
y
z
M x y z
+
+
+
+
+
-
\in \Pi \Leftrightarrow 
=
+
+
++
,即 
x
y
z
+
+
-
=
,或者
x
y
z
+
-
-
=
. 
\textcolor{lyred}{补充练习 }
1.求与球面
x
y
z
x
z
+
+
-
+
+
=
切于
(
)
2,
2, 4
M
-
的平面方程. 
解.(
)
(
)
x
y
z
-
+
+
+
=
,球心为
(
)
1,0, 3
A
-
,
(
)
1,
2, 1
n
AM
=
=
-
\xi \xi \xi \xi \rho 
\rho 
,故 
(
)
(
)
(
)
x
y
z
\cdots 
-
+
\cdots 
-
-\cdots 
+
=
,即
x
y
z
+
-
-
=
. 
2.求与原点的距离为6 ,三个截距之比为1:3: 2 的平面方程. 
解.设
x
y
z
x
y
z
k
k
k
k
+
+
=\Rightarrow 
+
+
-
=
,由
k
k
-
=
\Rightarrow 
=\pm 
+
+
,故 
x
y
z
+
+
\pm 
=
. 
3.求圆
x
y
z
x
y
z
x
y
z

+
+
-
+
-
=

+
-
=

的面积. 
解.(
)
(
)
(
)
x
y
z
-
+
+
+
-
=
,球心(
)
1, 2,3
-
到2
x
y
z
+
-
=的距离
d =,故 
上述圆的半径
r =
-
=
,于是
A
\pi 
=
. 
4.已知
(
)
1,0,2
A
,
(
)
1,1,1
B -
,
(
)
3,0,4
C
,求平行于三点所在平面,并且与该平面的 
距离为2 的平面方程. 
解.三点所在平面为
x
y
z
+
-
+=
,设所求平面为
x
y
z
D
+
-
+
=
,则 
2 3
1 1 1
D
D
-
=
\Rightarrow 
=\pm 
++
,故
2 3
x
y
z
+
-
+\pm 
=
. 
\vspace{6pt}
{\large \bfseries \textcolor{lyred}{第8.4 节 空间直线及其方程}}

\vspace{4pt}
{\bfseries \textcolor{lyred}{一.空间直线的一般方程}}

方程组
A x
B y
C z
D
A x
B y
C z
D
+
+
+
=


+
+
+
=

称为空间直线的\textcolor{lyred}{一般方程}. 
\textcolor{lyred}{例}.求过点(
)
3,2,5
-
,且与直线
x
z
x
y
z
-
=


-
-
=

平行的直线方程. 
解.过(
)
3,2,5
-
,平行于
x
z
-
=
的平面为
x
z
-
+
=
;过(
)
3,2,5
-
,平行于 
x
y
z
-
-
=的平面为2
x
y
z
-
-
+
=
,故所求直线为
x
z
x
y
z
-
+
=


-
-
+
=

. 
\vspace{4pt}
{\bfseries \textcolor{lyred}{二.平面束方程}}

\textcolor{lyred}{定理}.设平面\Pi 过
:
A x
B y
C z
D
L
A x
B y
C z
D
+
+
+
=


+
+
+
=

,则存在,
\lambda \mu \in \mathbb{R},使得\Pi 可表示为 
(
)
(
)
A x
B y
C z
D
A x
B y
C z
D
\mu 
\lambda 
+
+
+
+
+
+
+
=
. 
\textcolor{lyred}{推论}.若\Pi 不是平面
A x
B y
C z
D
+
+
+
=
(
\mu \neq 
),则它的方程可以表示为 
(
)
A x
B y
C z
D
A x
B y
C z
D
\lambda 
+
+
+
+
+
+
+
=
. 
\textcolor{lyred}{例}.求过点(
)
1,2,1 和直线
x
y
z
x
z
-
-
=

-
=

的平面方程. 
解.设
(
)
x
y
z
x
z
\lambda 
-
-
+
--
=
,代入(
)
1,2,1 ,得
\lambda =-,故
y
z
+
=
. 
\textcolor{lyred}{例}.求
:
x
y
z
L
x
y
z
+
-
-=

-
+
+=

在
1 :
x
y
z
\Pi 
+
+
=
上的投影直线L'的方程;并求过L', 
且与
\Pi 的夹角为6
\pi 的平面\Pi 的方程. 
解.设过L 且垂直于
\Pi 的平面为
(
)
x
y
z
x
y
z
\lambda 
+
--+
-
+
+
=
,由 
(
)(
)
,1
,
1,1,1
\lambda 
\lambda \lambda 
\lambda 
+
-
-
\cdots 
=
\Rightarrow 
=-,即
y
z
-
-=
,故
:
y
z
L
x
y
z
-
-=

'+
+
=

; 
设
(
)
:
y
z
x
y
z
\lambda 
\Pi 
--+
+
+
=
,则
(
)
,1
,
n
\lambda 
\lambda \lambda 
=
+
-
\rho 
,由 
(
)(
)
(
)(
)
,1
,
1,1,1
cos
,1
,
1,1,1
\lambda 
\lambda \lambda 
\lambda 
\pi 
\lambda 
\lambda 
\lambda \lambda 
\lambda 
+
-
\cdots 
=
\Rightarrow 
=
\Rightarrow 
=\pm 
+
-
+
\cdots 
,故\Pi 的方程为 
(
)
(
)
x
y
z
+
+
+
-
-=
或
(
)
(
)
x
y
z
-
+
-
+-
-
-=
. 
\vspace{4pt}
{\bfseries \textcolor{lyred}{三.空间直线的对称式方程}}

\textcolor{lyred}{定理}.设直线L 过点
(
)
,
,
M
x
y z
,且平行于非零向量
(
)
, ,
s
m n p
=
\rho 
,则它满足方程 
x
x
y
y
z
z
m
n
p
-
-
-
=
=
,称为L 的\textcolor{lyred}{对称式方程}. 
\textcolor{lyred}{注}.
(
)
, ,
s
m n p
=
\rho 
称为L 的\textcolor{lyred}{方向向量},m ,n , p 称为L 的\textcolor{lyred}{方向数}. 
\textcolor{lyred}{推论}.过两点
(
)
,
,
i
i
i
i
M
x y z
的直线方程为
x
x
y
y
z
z
x
x
y
y
z
z
-
-
-
=
=
-
-
-
. 
\textcolor{lyred}{例}.求过点(
)
3,2,5
-
,且平行于直线
x
y
z
-
+
=
=
-
的直线方程. 
解.
(
)
1,2,0
s =-
\rho 
,故
x
y
z
+
-
-
=
=
-
. 
\textcolor{lyred}{例}.求过点(
)
1, 2,4
-
,且垂直于平面2
x
y
z
-
+
-
=
的直线方程. 
解.
(
)
2, 3,1
s =
-
\rho 
,故
x
y
z
-
+
-
=
=
-
. 
\textcolor{lyred}{例}.求过点(
)
1,0, 3
-
,且与平面
x
z
-
=
和2
x
y
z
-
-
=平行的直线方程. 
解.
(
)(
)
(
)
1,0, 4
2, 1, 5
4, 3, 1
s =
-
\times 
--
=-
-
-
\rho 
,故
x
y
z
-
+
=
=
. 
\textcolor{lyred}{例}.求直线
: 2
x
y
z
L
x
y
z
+
+
+=


-
+
+
=

的对称式方程. 
解.令
x =,由
y
z
y
y
z
+
=-

\Rightarrow 
=
-
=

,
z =-,故
(
)
1,0, 2
M
L
-
\in 
; 
再取
(
)(
)
(
)
1,1,1
2, 1,3
4, 1, 3
s =
\times 
-
=
--
\rho 
,得
:
x
y
z
L
-
+
=
=
-
-
. 
\vspace{4pt}
{\bfseries \textcolor{lyred}{四.空间直线的参数方程}}

\textcolor{lyred}{定理}.过点(
)
,
,
x
y z
,且平行于非零向量(
)
, ,
m n p 的直线可表示为
x
mt
x
y
nt
y
z
pt
z
=
+


=
+

=
+

. 
\textcolor{lyred}{例}.求直线
x
y
z
-
-
-
=
=
与平面2
x
y
z
+
+
-
=
的交点. 
解.设交点为
(
)
2,
3,2
Q t
t
t
+
+
+
,代入2
x
y
z
+
+
-
=
,得
t =-,故
(
)
1,2,2
Q
. 
\textcolor{lyred}{例}.求点
(
)
1,2, 2
P -
-
关于平面
:
x
y
z
\Pi 
+
-
+=
的对称点
1P 的坐标. 
解.过P ,垂直于\Pi 的直线
:
x
y
z
L
+
-
+
=
=-
,设
(
)
1,2
2,
L
Q t
t
t
\cap \Pi =
-
+
--
, 
代入
x
y
z
+
-+=,得
t =-,故
(
)
2,0, 1
Q -
-
,于是
(
)
3, 2,0
P -
-
. 
\textcolor{lyred}{例}.设一束光线由(
)
3, 5,2
P
-
出发,沿直线9
x
y
z
+
-
=

=

射向
x
y
z
-
-
+
=
, 
求反射光线的方程. 
解.光线与平面的交点为
(
)
1,4,2
Q
,而P 关于平面的对称点为
(
)
1 1, 3,8
P
-
,故 
反射光线即为直线
:
x
y
z
PQ
-
-
-
=
=-
. 
\textcolor{lyred}{例}.求过点(
)
2, 1,3
P
-
,且与直线
: 3
x
y
z
L
+
-
=
=
垂直相交的直线方程. 
解.设两直线交点为
(
)
3 ,5
7,2
Q
t
t
t
-
+
,则
(
)
2,5
6,2
PQ
t
t
t
=
-
-
-
\xi \xi \xi \rho 
,由 
(
)
3,5,2
PQ
t
\cdots 
=
\Rightarrow =
\xi \xi \xi \rho 
,得交点
(
)
3, 2,4
Q
-
,故
x
y
z
-
+
-
=
=
-
. 
\textcolor{lyred}{例}.求过点
(
)
1,2,3
P
,且与直线
:
x
y
z
L
x
y
z
+
-
+=

-
+
-=

垂直相交的直线方程. 
解.过P 且垂直于L 的平面方程为
x
y
z
-
-
+
=
,由 
x
y
z
x
y
z
x
y
z
+
-
+=

-
+
-=

-
-
+
=

,得交点
(
)
1,2,2
Q -
,故
x
y
z
-
-
-
=
=
-
-
. 
\textcolor{lyred}{例}.求过(
)
2,1,0
P
,平行于
:3
x
y
z
\Pi 
-
-
=,与
x
y
z
-
+
=
=-
相交的直线方程. 
解.设两直线交点为
(
)
2 ,3
1,
Q
t
t
t
+
--
,则
(
)
2,3 ,
PQ
t
t
t
=
-
--
\xi \xi \xi \rho 
,由 
(
)
(
)
3, 2, 2
2,6, 3
PQ
t
PQ
\cdots 
-
-
=
\Rightarrow =
\Rightarrow 
=
-
\xi \xi \xi \rho 
\xi \xi \xi \rho 
,故
x
y
z
-
-
=
=-. 
\textcolor{lyred}{例}.求过
(
)
1,0,4
M -
,与
1 : 1
x
y
z
L
=
=
,
:
x
y
z
L
-
-
-
=
=
均相交的直线方程. 
解一.过M 与
1L 的平面为8
x
y
z
-
+
=
(过L ),过M 与
L 的平面为 
x
y
z
-
-
+
=
(也过L ),故所求直线为8
x
y
z
x
y
z
-
+
=


-
-
+
=

. 
解二.
(
)
0,0,0
M
L
\in 
,
(
)
1,2,3
M
L
\in 
,设
(
)
, ,
s
m n p
=
\rho 
,则由 
,
,
,
,
m
n
p
MM s s
MM
s s
m
n
p
m
n
p
-
+
=





=
=
\Rightarrow 
\Rightarrow 
=
=-





-
-
=

\xi \xi \xi \xi \xi \rho 
\xi \xi \xi \xi \xi \rho 
\rho \rho 
\rho \rho 
,故 
所求直线为
x
y
z
+
-
=
=
-
-
. 
\vspace{4pt}
{\bfseries \textcolor{lyred}{五.两直线的夹角}}

两直线方向向量的夹角
\pi 
\theta 
\theta 


\le 
\le 




称为两直线的\textcolor{lyred}{夹角}.因此,若 
: x
x
y
y
z
z
L
m
n
p
-
-
-
=
=
,
: x
x
y
y
z
z
L
m
n
p
-
-
-
=
=
,则夹角余弦为 
cos
s s
m m
n n
p p
s s
m
n
p
m
n
p
\theta 
\cdots 
+
+
=
=
+
+
+
+
\rho \rho 
\rho \rho 
. 
\textcolor{lyred}{例}.求直线
:
x
y
z
L
-
-
-
=
=
-
与
: 2
x
y
L
y
z
-
=


+
=

的夹角余弦. 
解.
(
)
1, 2,1
s =
-
\rho 
,
(
)(
)
(
)
1, 1,0
0,2,1
1, 1,2
s =
-
\times 
=--
\rho 
,故
cos
s s
s s
\theta 
\cdots 
=
=
\rho \rho 
\rho \rho 
. 
\vspace{4pt}
{\bfseries \textcolor{lyred}{六.直线与平面的夹角}}

\textcolor{lyred}{定义}.直线
x
x
y
y
z
z
m
n
p
-
-
-
=
=
与它在平面
Ax
By
Cz
D
+
+
+
=
上的投影直线 
之间的夹角
\pi 
\varphi 
\varphi 


\le 
\le 




,称为该直线与平面的\textcolor{lyred}{夹角},因此 
sin
s n
Am
Bn
Cp
s n
A
B
C
m
n
p
\varphi 
\cdots 
+
+
=
=
+
+
\cdots 
+
+
\rho \rho 
\rho \rho 
. 
\vspace{4pt}
{\bfseries \textcolor{lyred}{七.点到直线的距离}}

\textcolor{lyred}{例}.求点
(
)
1,2,3
M
到直线
: 1
x
y
z
L
-
-
=
=
-
的距离. 
解.设M 到L 的垂足为
(
)
, 3
4,2
N t
t
t
-
+
+
,则
(
)
1, 3
2,2
MN
t
t
t
=
-
-
+
\xi \xi \xi \xi \rho 
,由 
(
)
1, 3,2
MN
t
\cdots 
-
=
\Rightarrow =
\xi \xi \xi \xi \rho 
,故
1 5
,
,4
2 2
N 





,于是
d
MN
=
=
. 
\textcolor{lyred}{注}.也可先求过M ,且垂直于L 的平面\Pi ,然后再求出N
L
=
\cap \Pi . 
\textcolor{lyred}{定理}.点M 到直线L 的距离
MN
s
d
s
\times 
=
\xi \xi \xi \xi \rho 
\rho 
\rho 
,其中N
L
\in 
. 
\textcolor{lyred}{定理}.异面直线
1L ,
2L 的距离
(
)
M M
s
s
d
s
s
\cdots 
\times 
=
\times 
\xi \xi \xi \xi \xi \xi \rho 
\rho 
\rho 
\rho 
\rho 
,其中
M
L
\in 
,
M
L
\in 
. 
\textcolor{lyred}{注}.直线
1L 与
2L 共面
(
)
M M
s
s
\Leftrightarrow 
\cdots 
\times 
=
\xi \xi \xi \xi \xi \xi \rho 
\rho 
\rho 
. 
\textcolor{lyred}{例}.设直线
1 : 1
x
y
z
L
=
=
,
:
x
y
z
L
-
+
-
=
=
,(1)证明它们为异面直线; 
(2)求公垂线的方程. 
(1)证.
(
)
0,0,0
M
L
\in 
,
(
)
1, 1,2
M
L
-
\in 
,
(
)
M M
s
s
\cdots 
\times 
=-\neq 
\xi \xi \xi \xi \xi \xi \rho 
\rho 
\rho 
,证毕. 
(2)解.
(
)(
)
(
)
1,2,3
1,1,1
1,2, 1
s =
\times 
=-
-
\rho 
,过
1L 平行于s\rho 的平面为 
x
y
z
+
-
=
,过
2L 平行于s\rho 的平面为
x
z
-
+=
,故4
x
y
z
x
z
+
-
=

-
+=

. 
\textcolor{lyred}{补充练习 }
1.证明:
x
y
z
-
-
+
=
=-
与
x
y
z
-
-
=
=-
相交,并求它们确定的平面方程. 
解.
x
z
x
y
y
x
y
z
z
-
+


=

=
-


=
\Rightarrow 
=




-
-


=
=
=
-


,故有交点
1,2, 2






,即相交; 
(
)(
)
(
)
4,0, 3
2,2, 1
6, 2,8
n =
-
\times 
-
=
-
\rho 
,故
(
)
(
)
(
)
x
y
z
\cdots 
-
-\cdots 
-
+\cdots 
-
=
. 
2.求过
: 4
x
y
L
x
y
z
+
=


+
+
=

,且与球面
x
y
z
+
+
=
相切的平面方程. 
解.设
(
)
x
y
z
x
y
\lambda 
+
+
-
+
+
=
,则
(
)
(
)
\lambda 
\lambda 
\lambda 
=
\Rightarrow 
=-
+
+
+
+
, 
故
z =
. 
3.求过(
)
1,2,1
P
,且与
x
y
z
-
+
=
=
垂直,与2
x
y
z
=
=-相交的直线方程. 
解.设两直线的交点为
(
)
2 , ,
Q
t t
t
-
,由
(
)
(
)
1,
2,
3,2,1
PQ
t
t
t
=
-
-
--
\bot
\xi \xi \xi \rho 
,得
t =
, 
即
,
,
PQ


=
-
-




\xi \xi \xi \rho 
,于是
x
y
z
-
-
-
=
=
-
-
. 
4.求
:
x
y
z
\Pi 
+
+
=上与
:
x
y
z
L
-
+
=
=-
垂直相交的直线方程. 
解.设L 与\Pi 的交点为
(
)
1, ,
M t
t
t
+
--
,代入
x
y
z
+
+
=,得
t =
,故
(
)
3,2, 4
M
-
, 
而
(
)(
)
(
)
1,1, 1
1,1,1
2, 2,0
s
s
n
=
\times 
=
-
\times 
=
-
\rho 
\rho 
\rho 
,于是
x
y
z
-
-
+
=
=
-
. 
5.设直线
1L 经过点
(
)
2, 3,5
M
-
,并且与x 轴以及y 轴的夹角均为3
\pi ,求
1L 与直线 
2 :
L
x
y
z
=
=
的距离. 
解.
cos
cos
cos
cos
\pi 
\alpha 
\beta 
\gamma 
=
=
=
\Rightarrow 
=
-
-
=
,故
(
)
1,1,
s =
\rho 
, 
取
(
)
2, 3,5
M
L
-
\in 
,
(
)
0,0,0
N
L
\in 
,则
(
)
NM
s
s
d
s
s
\cdots 
\times 
=
=
\times 
\xi \xi \xi \xi \rho 
\rho 
\rho 
\rho 
\rho 
. 
6.在平面
x
y
z
-
+
=上求一点M ,使得它到
(
)
1,0, 1
P
-
与
(
)
2,1,2
Q
的距离之和 
最小. 
解.(1)设
(
)
, ,
F x y z
x
y
z
=
-
+
-,则
(
)
1,0, 1
F
-
=-,
(
)
2,1,2
F
=-,故两点在 
该平面的同一侧; 
(2)求得P 关于平面的对称点为
(
)
1 3, 4,3
P
-
; 
(3)直线
1PQ 与平面的交点
7 30
,
,
13 13
M 

-




即为所求. 
\vspace{6pt}
{\large \bfseries \textcolor{lyred}{第8.5 节 曲面及其方程}}

\vspace{4pt}
{\bfseries \textcolor{lyred}{一.曲面的一般方程}}

设有曲面S ,若(
)
(
)
, ,
, ,
x y z
S
F x y z
\in 
\Leftrightarrow 
=
,则称方程
(
)
, ,
F x y z =
为S 的\textcolor{lyred}{方程}, 
而S 称为方程
(
)
, ,
F x y z =
的\textcolor{lyred}{图形}. 
\textcolor{lyred}{例}.设点(
)
1,2,3
A
, (
)
2, 1,4
B
-
,求线段AB 的垂直平分面方程. 
解.
(
)
, ,
M x y z
S
\forall 
\in 
,则AM
BM
=
,即(
)
(
)
(
)
x
y
z
-
+
-
+
-
= 
(
)
(
)
(
)
x
y
z
-
+
+
+
-
,因此2
x
y
z
-
+
-
=
. 
\textcolor{lyred}{例}.求球心在
(
)
,
,
M
x
y z
,半径为R 的球面方程. 
解.
(
)
, ,
M x y z
S
\forall 
\in 
,则
MM
R
=
,即(
)
(
)
(
)
x
x
y
y
z
z
R
-
+
-
+
-
=
. 
\textcolor{lyred}{例}.方程
x
y
z
x
y
+
+
-
+
=
表示怎样的曲面? 
解.(
)
(
)
x
y
z
-
+
+
+
=
,表示球心在(
)
1, 2,0
-
,半径5 的球面. 
\vspace{4pt}
{\bfseries \textcolor{lyred}{二.旋转曲面}}

\textcolor{lyred}{定义}.一条曲线绕定直线旋转一周,所形成的曲面称为\textcolor{lyred}{旋转曲面},定直线称为它的 
\textcolor{lyred}{旋转轴},动曲线称为它的\textcolor{lyred}{母线}. 
\textcolor{lyred}{例}.设
(
)
:
,
C f y z =为yOz 面上的平面曲线,把它绕z 轴旋转一周,得到旋转面S , 
求S 的方程. 
解.
(
)
, ,
M x y z
S
\forall 
\in 
,
(
)
0,
,
M
y z
C
\exists 
\in 
,使得它们位于同一个旋转圆上,于是 
y
x
y
z
z

=
+

=

,由
(
)
,
f
y z
=
,得
(
)
:
,
S
f
x
y
z
\pm 
+
=
. 
\textcolor{lyred}{注}.若上述曲线绕y 轴旋转,则所得曲面方程为(
)
,
f
y
x
z
\pm 
+
=
. 
类似地,可以得到xOy 面( xOz 面)上的曲线绕,x y 轴( ,x z 轴)旋转所得到的曲面 
方程. 
\textcolor{lyred}{例}.设
:
C z
ky
=
为yOz 面上的直线,它绕z 轴旋转一周,得到曲面S ,它的方程为 
z
k
x
y
=\pm 
+
,即
(
)
z
k
x
y
=
+
,称为\textcolor{lyred}{圆锥面}. 
\textcolor{lyred}{例}.设
:
y
C z
a
=
为yOz 面上的抛物线,绕z 轴旋转一周,得到曲面S ,它的方程为 
x
y
z
a
+
=
,称为\textcolor{lyred}{旋转抛物面}. 
\textcolor{lyred}{例}.设
:
x
z
C a
c
-
=为xOz 面上的双曲线,求:(1)它绕z 轴旋转一周形成的曲面 
方程;(2)它绕x 轴旋转一周形成的曲面方程. 
解.(1)绕z 轴旋转:
x
y
z
a
c
+
-
=,\textcolor{lyred}{单叶旋转双曲面}; 
(2)绕x 轴旋转:
x
y
z
c
a
+
-
=,\textcolor{lyred}{双叶旋转双曲面}. 
\textcolor{lyred}{例}.设
:
x
L
y
t
z
t
=


=

=

,求L 绕z 轴旋转所得曲面的方程. 
解.
(
)
, ,
M x y z
S
\forall 
\in 
,
(
)
1 1, ,2
M
t
t
L
\exists 
\in 
,使得它们在一个旋转圆上,于是 
x
y
t
z
z
x
y
x
y
z
t

+
=+


\Rightarrow 
+
=+
\Rightarrow 
+
-
=



=



,单叶旋转双曲面. 
\vspace{4pt}
{\bfseries \textcolor{lyred}{三.柱面}}

\textcolor{lyred}{定义}.设动直线L 沿曲线\Gamma 平行移动,则它所形成的轨迹称为\textcolor{lyred}{柱面},称\Gamma 为\textcolor{lyred}{准线}, 
L 为\textcolor{lyred}{母线}. 
\textcolor{lyred}{例}.求平行于z 轴,以xOy 面上
(
)
:
,
C
f x y =
为准线的柱面方程. 
解.
(
)
, ,
M x y z
S
\forall 
\in 
,
(
)
, ,0
M
x y
C
'
\in 
,故
(
)
,
f x y =
. 
\textcolor{lyred}{注}.
(
)
,
f x z =
表示平行于y 轴的柱面;
(
)
,
f
y z =
表示平行于x 轴的柱面. 
\textcolor{lyred}{例}.求平行于
(
)
3,4,1
a =-
\rho 
,以xOy 面上的曲线
:
x
y
C a
b
+
=为准线的柱面方程. 
解.
(
)
, ,
M x y z
S
\forall 
\in 
,
(
)
,
,0
M
x y
C
'
\exists 
\in 
,使得
x
x
y
y
z
-
-
-
=
=
-
,即 
x
x
z
y
y
z
=
+


=
-

,故
(
)
(
)
:
x
z
y
z
S
a
b
+
-
+
=. 
\vspace{4pt}
{\bfseries \textcolor{lyred}{四.二次曲面}}

三元二次方程所表示的曲面称为\textcolor{lyred}{二次曲面},共有九大类. 
\textcolor{lyred}{椭圆柱面}:
x
y
a
b
+
=;\textcolor{lyred}{双曲柱面}:
x
y
a
b
-
=;\textcolor{lyred}{抛物柱面}:
2x
ay
=
; 
\textcolor{lyred}{椭圆锥面}:
x
y
z
a
b
+
=
;\textcolor{lyred}{椭球面}:
x
y
z
a
b
c
+
+
=; 
\textcolor{lyred}{单叶双曲面}:
x
y
z
a
b
c
+
-
=;\textcolor{lyred}{双叶双曲面}:
x
y
z
a
b
c
-
-
=; 
\textcolor{lyred}{椭圆抛物面}:
x
y
z
a
b
+
=
;\textcolor{lyred}{双曲抛物面}:
x
y
z
a
b
-
=
,又称\textcolor{lyred}{马鞍面}. 
\textcolor{lyred}{补充练习 }
1.求经过点(
)
1,0, 2
-
的平面与球面
(
)
(
)
x
y
z
+
-
+
+
=
的交线的最小长度. 
解.(
)
1,0, 2
-
到球心的距离
d =
,故最小交线的半径
r =
-
=
,而
l
\pi 
=
. 
2.求顶点在原点O ,轴平行于
(
)
, ,
a
m n p
=
\rho 
,半顶角为
\pi 
\theta 
\theta 


<
<




的圆锥面方程. 
解.
(
)
, ,
M x y z
S
\forall 
\in 
,则OM
\xi \xi \xi \xi \rho 
与a\rho 的夹角为\theta ,或者\pi 
\theta 
-
,故 
(
)
(
)(
)
cos
cos
mx
ny
pz
OM a
x
y
z
m
n
p
OM a
\theta 
\theta 
+
+
\cdots 
=\pm 
\Rightarrow 
=
+
+
+
+
\xi \xi \xi \xi \rho \rho 
\xi \xi \xi \xi \rho \rho 
,即 
(
)
(
)(
)
cos
mx
ny
pz
m
n
p
x
y
z
\theta 
+
+
=
+
+
+
+
. 
3.求半径为R ,对称轴过原点且平行于
(
)
, ,
a
m n p
=
\rho 
的圆柱面方程. 
解.
(
)
, ,
M x y z
S
\forall 
\in 
,则(
)
Prja OM
R
OM
+
=
\rho \xi \xi \xi \xi \rho 
\xi \xi \xi \xi \rho 
,故 
(
)
mx
ny
pz
R
x
y
z
m
n
p
+
+
+
=
+
+
+
+
. 
4.设
:
x
t
L
y
t
z
t
=


=
-

=-+

,求L 绕x 轴旋转所得曲面的方程. 
解.
(
)
, ,
M x y z
S
\forall 
\in 
,
(
)
1 2 ,2
2,
M
t
t
t
L
\exists 
-
-+
\in 
,它们在一个旋转圆上,于是 
(
)
(
)
y
z
t
t
x
t

+
=
-
+-+
=

,消去t ,得
(
)
x
y
z
x


+
=
-
+
-




,即 
(
)
x
y
z
-
-
-
=
,圆锥面. 
\vspace{6pt}
{\large \bfseries \textcolor{lyred}{第8.6 节 空间曲线及其方程}}

\vspace{4pt}
{\bfseries \textcolor{lyred}{一.空间曲线的一般方程}}

空间曲线\Gamma 可视为两个曲面的交,若已知两个曲面的方程分别为
(
)
, ,
F x y z =
和 
(
)
, ,
G x y z =
,则(
)
(
)
(
)
, ,
, ,
, ,
F x y z
x y z
G x y z
=

\in \Gamma \Leftrightarrow 
=

,此方程组称为\Gamma 的\textcolor{lyred}{一般方程}. 
\vspace{4pt}
{\bfseries \textcolor{lyred}{二.空间曲线的参数方程}}

空间曲线也可表示为\Gamma :
()
()
()
x
x t
y
y t
z
z t
=


=


=

,t
I
\in 
,当t 在I 上变动时,对应的(
)
, ,
x y z 走过 
曲线上的所有点,称为空间曲线的\textcolor{lyred}{参数方程}. 
\textcolor{lyred}{例}.设M 位于柱面
x
y
a
+
=
上,以角速度\omega 绕z 轴旋转,以线速度v 沿z 轴上升, 
它的轨迹称为\textcolor{lyred}{圆柱螺旋线},求它的参数方程. 
解.取时间t 为参数,设
t =
时M 在点
(
)
,0,0
A a
,则轨迹的参数方程为 
(
)
cos
sin
x
a
t
y
a
t t
z
vt
\omega 
\omega 
=


=
\ge 

=

,也可以取
t
\theta 
\omega 
=
为参数,则
(
)
cos
:
sin
x
a
y
a
z
b
\theta 
\theta \theta 
\theta 
=


\Gamma 
=
\ge 

=

. 
\textcolor{lyred}{例}.设
:
x
t
y
t
z
t
=


\Gamma 
=

=

,求\Gamma 绕z 轴旋转所得曲面的方程. 
解.
(
)
, ,
M x y z
S
\forall 
\in 
,
(
)
,
,
M
t t
t
\exists 
\in \Gamma ,使得它们在一个旋转圆上,于是 
()
x
y
t
t
z
t

+
=
+

=

,消去t ,得到
x
y
z
z
+
=
+
. 
\vspace{4pt}
{\bfseries \textcolor{lyred}{三.曲面的参数方程}}

曲面也可以表示为
(
)
(
)
(
)
,
:
,
,
x
x s t
S
y
y s t
z
z s t
=


=


=

,(
)
,s t
D
\in 
,称为曲面的\textcolor{lyred}{参数方程}. 
\textcolor{lyred}{例}.设
()
()
()
:
x
x t
y
y t
z
z t
=


\Gamma 
=


=

,
[
]
,
t
\alpha \beta 
\in 
,求\Gamma 绕z 轴旋转所得曲面的参数方程. 
解.
(
)
, ,
M x y z
S
\forall 
\in 
,
()
()
()
(
)
,
,
M
x t
y t
z t
\exists 
\in \Gamma ,它们在一个旋转圆上,故 
()
()
()
()
()
()
()
()
cos
sin
x
x
t
y
t
x
y
x
t
y
t
y
x
t
y
t
z
z t
z
z t
\theta 
\theta 
=
+


+
=
+


\Rightarrow 
=
+


=


=

,其中
[
]
,
t
\alpha \beta 
\in 
,
[
]
0,2
\theta 
\pi 
\in 
. 
\textcolor{lyred}{例}.求半径为a 的球面
:
S x
y
z
a
+
+
=
的参数方程. 
解.取yOz 面上半圆
:
sin
cos
x
y
a
z
a
\varphi 
\varphi 
=


\Gamma 
=

=

,
[
]
0,
\varphi 
\pi 
\in 
,它绕z 轴旋转一周形成球面S ,故 
sin
cos
:
sin
sin
cos
x
a
S
y
a
z
a
\varphi 
\theta 
\varphi 
\theta 
\varphi 
=


=

=

,
[
]
0,
\varphi 
\pi 
\in 
,
[
]
0,2
\theta 
\pi 
\in 
. 
\vspace{4pt}
{\bfseries \textcolor{lyred}{四.空间曲线在坐标面上的投影}}

\textcolor{lyred}{定义}.以空间曲线\Gamma 为准线,平行于z 轴的柱面S ,称为\Gamma 关于xOy 面的\textcolor{lyred}{投影柱面}, 
它与xOy 面的交线'
\Gamma 称为\Gamma 在xOy 面上的\textcolor{lyred}{投影曲线}. 
设
(
)
(
)
, ,
:
, ,
F x y z
G x y z
=

\Gamma 
=

,则
(
)
, ,
M x y z
S
\forall 
\in 
,存在
1z ,使得
(
)
, ,
M
x y z \in \Gamma ,由 
(
)
(
)
, ,
, ,
F x y z
G x y z
=

=

,消去
1z ,得
(
)
:
,
S H x y =
,于是
(
)
,
:
H x y
z
=

'
\Gamma 
=

. 
\textcolor{lyred}{例}.设
(
)
(
)
:
x
y
z
x
y
z

+
+
=

\Gamma 
+
-
+
-
=

,求\Gamma 在xOy 面上的投影曲线的方程. 
解.
(
)
y
z
z
y
x
y
y
x
y
y
+
=\Rightarrow 
=-
\Rightarrow 
+
+
-
=\Rightarrow 
+
-
=
,即 
x
y


+
-
=




. 
\textcolor{lyred}{例}.求
z
x
y
=
-
-
和
(
)
z
x
y
=
+
所围立体在xOy 面上的投影区域. 
解.由
(
)
x
y
x
y
x
y
-
-
=
+
\Rightarrow 
+
=,故投影区域为
x
y
+
\le . 
\vspace{4pt}
{\bfseries \textcolor{lyred}{五.由空间曲线的一般方程建立参数方程}}

\textcolor{lyred}{例}.设
:
x
y
x
y
z

+
=
\Gamma -
+
=

,求\Gamma 的参数方程. 
解.
cos
:
sin
cos
sin
x
t
y
t
z
t
t
=


\Gamma 
=

=
-
+

. 
\textcolor{lyred}{例}.设
:
z
x
y
x
y
x
=
-
-
\Gamma 
+
=

,求\Gamma 的参数方程.\textcolor{lyred}{ }
解.
(
)
x
y
x
x
y
+
=
\Rightarrow 
-
+
=,于是
1 cos
:
sin
2cos
2sin 2
x
t
y
t
t
z
t

=+

\Gamma 
=


=
-
=

,
[
)
0,2
t
\pi 
\in 
. 
\textcolor{lyred}{补充练习 }
1.设
(
)
(
)
:
z
x
y
z
x
y
=
-
-

\Gamma =
-
+
-

,求\Gamma 在三个坐标面上的投影曲线的方程. 
解.消去z ,得
x
y
x
y
+
=
+
,故在xOy 面上的投影为
x
y
x
y
z

+
=
+
=

; 
消去y ,得
(
)
z
x
x
z
x
xz
z
x
z
=
-
-
-
-
\Rightarrow 
+
+
-
-
+
=
,故 
在xOz 面上的投影曲线为
x
xz
z
x
z
y

+
+
-
-
+
=

=

; 
由对称性,在yOz 面上的投影曲线为
y
yz
z
y
z
x

+
+
-
-
+
=
=

. 
2.求
z
x
y
=
+
与
z
x
y
=
-
-
所围立体在xOy 面上的投影区域. 
解.由
z
x
y
x
y
z
x
y
=
+
\Rightarrow 
+
=
=
-
-

,故投影区域为
x
y
+
\le 
. 